\subsection{PRODUCT OF TWO EQUIVALENCE RELATIONS}

Let \(R\{x, y\}\) and \(R'\{x', y'\}\) be two equivalence relations. Let \(S\{u, v\}\) denote the relation

\((\exists x)(\exists y)(\exists x^{\prime})(\exists y^{\prime})(u = (x,x^{\prime})\) and \(v = (y,y^{\prime})\) and \(\mathbf{R}\{x,y\}\) and \(\mathbf{R}'\{x',y'\}\) ;

it is immediately verified that \(S\{u, v\}\) is an equivalence relation, called the product of R and \(R'\) , and denoted by \(R \times R'\) . Suppose that R is an equivalence relation on a set E and that \(R'\) is an equivalence relation on a set \(E'\) . Then the relation \(S\{u, u\}\) is equivalent to

\[
(\exists x) (\exists x ^ {\prime}) (u = (x, x ^ {\prime}) \text { and } R \{x, x \} \text { and } R ^ {\prime} \{x ^ {\prime}, x ^ {\prime} \})
\]

i.e., to \((\exists x)(\exists x')(u = (x, x')\) and \(x \in \mathbf{E}\) and \(x' \in \mathbf{E}'\) ), hence to \(u \in \mathbf{E} \times \mathbf{E}'\) . It follows that \(\mathbf{R} \times \mathbf{R}'\) is an equivalence relation on \(\mathbf{E} \times \mathbf{E}'\) . If

\[
u = (x, x ^ {\prime})
\]

is an element of \(\mathbf{E} \times \mathbf{E}'\) , the relation \(S\{u, v\}\) is equivalent to

\[
(\exists y) (\exists y ^ {\prime}) (v = (y, y ^ {\prime}) \text {   and   } \mathrm{R} \{x, y \} \text {   and   } \mathrm{R} ^ {\prime} \{x ^ {\prime}, y ^ {\prime} \});
\]

if G and \(G'\) are the graphs of R and \(R'\) , respectively, this relation is in turn equivalent to \(v \in \mathbf{G}(x) \times \mathbf{G}(x')\) . Hence every equivalence class with respect to \(R \times R'\) is the product of an equivalence class with respect to R and an equivalence class with respect to \(R'\) , and conversely.

Let \(f\) and \(f'\) be the canonical mappings of \(\mathbf{E}\) onto \(\mathbf{E}/\mathbf{R}\) and \(\mathbf{E}'\) onto \(\mathbf{E}'/\mathbf{R}'\) , respectively, and let \(f \times f'\) be the canonical extension of \(f\) and \(f'\) to the product sets (§ 3, no. 9). Then \((f \times f')(x, x') = (f(x), f'(x'))\) for all \((x, x') \in \mathbf{E} \times \mathbf{E}'\) . The inverse image under \(f \times f'\) of an element \((u, u')\) of \((\mathbf{E}/\mathbf{R}) \times (\mathbf{E}'/\mathbf{R}')\) is just the product \(u \times u'\) of the equivalence class \(u\) with respect to \(\mathbf{R}\) and the equivalence class \(u'\) with respect to \(\mathbf{R}'\) . It follows that the equivalence relation associated with \(f \times f'\) is equivalent to \(\mathbf{R} \times \mathbf{R}'\) . The mapping \(f \times f'\) can therefore be written as \(h \circ g\) , where \(g\) is the canonical mapping of \(\mathbf{E} \times \mathbf{E}'\) onto

\[
(\mathbf {E} \times \mathbf {E} ^ {\prime}) / (\mathbf {R} \times \mathbf {R} ^ {\prime})
\]

and \(h\) is a one-to-one mapping of \((\mathbf{E} \times \mathbf{E}') / (\mathbf{R} \times \mathbf{R}')\) onto

\[
(\mathbf {E} / \mathbf {R}) \times (\mathbf {E} ^ {\prime} / \mathbf {R} ^ {\prime});
\]

this mapping h and its inverse are said to be canonical.

Remark. Let \(\mathrm{P}\{x, x'\}\) be a relation in which the letters \(y, y'\) do not occur. P is said to be compatible with the equivalence relations \(\mathrm{R}\{x, y\}\) and \(\mathrm{R}'\{x', y'\}\) (with respect to \(x\) and \(x'\) ) if the relation \((\mathrm{P}\{x, x'\}\) and \(\mathrm{R}\{x, y\}\) and \(\mathrm{R}'\{x', y'\}\) ) implies \(\mathrm{P}\{y, y'\}\) . Let \(\mathrm{Q}\{u\}\) be the relation \((\exists x)(\exists x')(u = (x, x')\) and \(\mathrm{P}\{x, x'\})\) ; then it comes to the same thing to say that \(\mathrm{Q}\{u\}\) is compatible (with respect to \(u\) ) with the equivalence relation \(\mathrm{S}\{u, v\}\) , the product of R and \(\mathrm{R}'\) .

